\bumppoints{50}

\question
One form of elephantiasis occurs when \textit{Brugia malayyi}, a thread-like worm, invades human lymph nodes. The worm blocks the flow of lymph throughout the body, causing massive swelling of the lower body, especially legs and genitals. As described here, this is an example of

\begin{choices}
\choice predation.
\correctchoice parasitism.
\choice mutualism.
\choice commensalism.
\choice symbiosis.
\end{choices}

\newpage

\question
A field containing 2000 kg of plant material can support about \rule{1in}{0.4pt} of tertiary consumers.

\begin{choices}
\correctchoice 2 kg %dontshuffle 
\choice 10\% %dontshuffle 
\choice 10 kg %dontshuffle 
\choice 20 kg %dontshuffle 
\choice 90\% %dontshuffle 
\choice 200 kg %dontshuffle 
\end{choices}

\question
Energy from the sun is converted into chemical energy for organisms by

\begin{choices}
\choice carnivores.
\correctchoice primary producers.
\choice herbivores.
\choice decomposers.
\choice primary consumers.
\end{choices}

\question
Competitive exclusion means that

\begin{choices}
\choice one population will always exclude another another population of the same species in a community. 
\choice two species will partition resources to eliminate competition. 
\correctchoice one species will always out complete another species that occupies the same niche.  
\choice one individual will always out complete another individual in a population. 
\choice a species will have a smaller realized niche than a fundamental niche. 
\end{choices}


\question
Organisms that break down organic material and make it available to other organisms are called

\begin{choices}
\choice scavengers. 
\correctchoice decomposers. 
\choice consumers. 
\choice producers.
\choice degraders.
\end{choices}

\question
The ``role'' or ``job'' of a species in a community, represented by all of the resources needed by the species, is called 

\begin{choices}
\choice a predator.
\correctchoice the niche.
\choice the resource.
\choice a predator or prey.
\choice a producer or consumer.
\end{choices}

\newpage

\question
Which block in the figure below shows the greatest species diversity? (Different shades, including white, represent different species. The small squares indicate relative abundance of individuals present.)

\begin{multicols}{3}
\begin{choices}
\choice \includegraphics{diversity_mcB} 
\choice \includegraphics{diversity_mcD} 
\choice \includegraphics{diversity_mcC} 
\choice \includegraphics{diversity_mcA} 
\correctchoice \includegraphics{diversity_mcE} 
\end{choices}
\end{multicols}

\question
Probably the most important factor(s) affecting the global distribution of ecosystems is (are)

\begin{choices}
\choice longitude and latitude.
\choice wind and ocean water current patterns.
\correctchoice climate.
\choice wind direction and rain shadows.
\choice day length and rainfall.
\end{choices}

\question
Colors and patterns that helps an animal blend in with its environment is called

\begin{choices}
\choice habitat mimicry. 
\choice aposomatic coloration. 
\correctchoice cryptic coloration. 
\choice environmental mimicry. 
\choice warning coloration 
\end{choices}

\question
An organism that only reproduces once is known as

\begin{choices}
\choice ectotherm. 
\choice iteroparous. 
\correctchoice semelparous. 
\choice perennial. 
\end{choices}

\question
What type of survivorship curve describes a species that has a high rate of juvenile mortality but low adult mortality?

\begin{choices}
\choice Type I curve. %dontshuffle 
\choice Type II curve. %dontshuffle 
\correctchoice Type III curve. %dontshuffle 
\choice Type IV curve %dontshuffle 
\end{choices}


\question
Based on the hypothetical food web below, one letter represents an organism that could be both a secondary and tertiary consumer? The arrows point in the direction of energy transfer from lower to higher trophic levels. More than one answer is possible. Circle only one answer.

\begin{multicols}{3} 
	\begin{choices}
\correctchoice A %dontshuffle 
\choice B %dontshuffle 
\choice C %dontshuffle 
\choice D %dontshuffle 
\correctchoice E %dontshuffle 
\end{choices}
\columnbreak
	\includegraphics{exam4_foodweb}
\columnbreak
\end{multicols}
\vspace{-1\baselineskip}

\question
Of the following choices, which is the biggest concern about climate change?

\begin{choices}
\choice More regions around the globe will experience stronger rain shadows. 
\choice The average global temperature will be warmer than ever recorded.  
\correctchoice The rate of temperature increase is faster than ever recorded. 
\choice The amount of solar radiation will change faster than plants can adapt. 
\choice Forest and grassland ecosystems will die due to less \ce{CO2} available for photosynthesis. 
\choice More regions around the globe will experience desert-like conditions. 
\end{choices}

\question
Early successional species add nutrients and deepen the soil, allowing intermediate successional species to colonize an area. This process is known as

\begin{choices}
\choice disturbance. 
\choice inhibition. 
\choice tolerance. 
\choice primary succession. 
\correctchoice facilitation. 
\choice secondary succession. 
\end{choices}

\question
The difference between communities and ecosystems is that

\begin{choices}
\choice communities are groups of individuals of the different species; ecosystems are groups of communities.
\choice communities include individuals of the same species; ecosystems include individuals of different species.
\choice communities include interactions between individuals and their physical environment; ecosystems do \emph{not} include such interactions.
\choice communities include individuals of the different species; ecosystems are individuals of the same species.
\correctchoice communities do \emph{not} include interactions between individuals and the physical environment; ecosystems \emph{do} include such interactions.
\end{choices}

\question
Which dispersion should be favored among plants that compete intraspecifically by casting   shade?

\begin{choices}
\choice Random. 
\choice Clumped. 
\choice Clustered. 
\correctchoice Evenly spaced. 
\end{choices}

\question
Which statement most accurately describes how nutrients and energy are used in ecosystems?

\begin{choices}
\choice Nutrients can be converted into energy; energy cannot be converted into nutrients. 
\correctchoice Energy flows through ecosystems; nutrients cycle within ecosystems. 
\choice Nutrients flow through ecosystems; energy cycles within ecosystems. 
\choice Nutrients are used in ecosystems; energy is not. 
\choice Energy can be converted into nutrients; nutrients cannot be converted into energy. 
\end{choices}

\question
White-breasted nuthatches and brown creepers both eat insects that hide in the furrows of bark in hardwood trees. The brown creeper searches for insects by hunting from the bottom of the tree trunk toward the top, whereas the white-breasted nuthatch searches from the top of the trunk down. Their feeding areas rarely overlap on the tree trunk. These hunting behaviors best illustrate which of the following ecological concepts?

\begin{choices}
\correctchoice resource partitioning
\choice keystone species
\choice competitive exclusion
\choice interspecific competition
\choice intraspecific competition
\end{choices}

\question
Fire suppression by humans

\begin{choices}
\choice is necessary for the protection of threatened and endangered forest species. 
\choice will result ultimately in sustainable production of increased amounts of forest products for human use. 
\correctchoice can change the species composition within biological communities. 
\choice is a management goal of conservation biologists to maintain the healthy condition of forest communities.
\choice will always result in an increase in species diversity in a given biome. 
\end{choices}

\question
Symbiosis is a long-term ecological interaction

\begin{choices}
\choice where both species benefit. 
\choice where one species benefits and another species is harmed 
\choice  where one species benefits and the other species neither benefits nor is harmed. 
\correctchoice between two or more species. 
\choice where two or more species partition resources. 
\end{choices}

\question
The energy that is \textit{not} transferred to other trophic levels or decomposers is lost

\begin{choices}
\choice in fecal matter.
\choice to primary consumers.
\correctchoice as metabolic heat.
\choice to secondary consumers.
\choice to consumers at any trophic level.
\end{choices}

\newpage

\question
Plants produce water vapor during photosynthesis, which goes from the plant to the atmosphere. The water vapor in the atmosphere eventually falls back to the earth as rain. This process would be part of

\begin{choices}
\choice energy flow.
\choice the nutrient cycle.
\choice the energy cycle.
\choice plant metabolism.
\correctchoice the water cycle.
\choice the entire ecosystem.
\end{choices}

\question
In a population growing according to the exponential growth model, population size is

\begin{choices}
\choice limited by density-independent factors. 
\choice limited by both density-dependent and density-independent factors. 
\correctchoice not limited. 
\choice limited by density-dependent factors. 
\end{choices}

\question
Which of the following statements about ecosystems is false?

\begin{choices}
\choice Water cycles within an ecosystem. 
\correctchoice About 10\% of the energy available at one trophic level is lost as it moves to the next higher trophic level. 
\choice Carnivores occupy higher trophic levels than herbivores. 
\choice Nutrients, like nitrogen and phosphorus, cycle within an ecosystem. 
\choice Energy flows through an ecosystem. 
\end{choices}

\question
According to this climograph, the temperate grassland 

\begin{multicols}{2}
	\begin{choices}
\choice never occurs with deserts. 
\choice has the greatest variation in precipitation of any ecosystem shown in the figure. 
\choice has an average annual rainfall of about 20--80 cm. 
\choice has an average annual rainfall of about 5--30 cm. 
\correctchoice has an average annual temperature range of about 4--30°C. 
\end{choices}
\columnbreak
	{\centering \includegraphics[width=0.9\columnwidth]{exam4_climograph1}\par}
\end{multicols}
